% Section 3: A Linear Type System for AI Accountability
% IEEE Computer — "Silent Decisions Are Type Errors"
% Authors: Soares et al., 2026

\section{A Linear Type System for AI Accountability}
\label{sec:type_system}

The legal accountability gap identified in Section~\ref{sec:intro}
stems from the mismatch between the \emph{sequential} structure of
runtime assertions and the \emph{concurrent, high-throughput} nature
of AI decision pipelines.  This section introduces the linear-type
framework that eliminates the gap at the language level.

% ----------------------------------------------------------------
\subsection{Regulatory Obligations as Linear Resources}
\label{sec:linear_resources}

Girard's linear logic~\cite{girard1987linear} treats propositions as
\emph{resources}: each assumption must be used exactly once---no
duplication (\emph{contraction}) and no silent discard
(\emph{weakening}).  This discipline maps directly onto regulatory
accountability obligations.

Under LGPD Art.~18\S2 and EU AI Act Art.~86, an AI system that
issues a consequential decision \emph{consumes} a regulatory
obligation: it must produce a corresponding explanation and preserve
an auditable evidence trail.  If the obligation could be duplicated,
the same evidence could be recycled across multiple decisions,
undermining non-repudiability.  If the obligation could be silently
discarded, decisions could escape accountability entirely---the
precise failure mode we call a \emph{silent decision}.

We model two regulatory obligations as first-class linear types:

\begin{description}
  \item[\texttt{EvidenceToken}]  A linear resource that encapsulates
    the BLAKE3 hash of the decision context (the raw input fed to the
    AI system at inference time).  Construction via
    \texttt{EvidenceToken::new(\&[u8])} binds the token irrevocably to
    a specific context; the hash is computed at construction and cannot
    be altered afterwards.  The type does not implement \texttt{Clone}
    or \texttt{Copy}, enforcing the ``used exactly once'' discipline at
    the type-system level.

  \item[\texttt{ComplianceToken}]  A linear resource that encodes the
    applicable regulatory framework: jurisdiction identifier, policy
    version, and contestability deadline in hours.  Like
    \texttt{EvidenceToken}, it is consumed by value when passed to the
    verdict constructor, ensuring that each regulatory obligation is
    bound to exactly one verdict.
\end{description}

The non-repudiability guarantee follows directly: because an
\texttt{EvidenceToken} cannot be cloned, the same evidence cannot be
presented for two different decisions.  Because it cannot be silently
dropped, a decision that lacks evidence cannot be compiled.

% ----------------------------------------------------------------
\subsection{The Atomic Binding Operator}
\label{sec:binding}

The central type law of this paper is:
\[
  V \multimap (E \otimes C)
\]
where $V$ denotes \texttt{Verdict}, $E$ denotes \texttt{EvidenceToken},
$C$ denotes \texttt{ComplianceToken}, $\multimap$ is the linear
function type (``linear implication''), and $\otimes$ is the
multiplicative conjunction (``tensor product'') of linear logic.

Reading right-to-left: to produce a $V$, one must simultaneously
consume exactly one $E$ \emph{and} exactly one $C$.  There is no
partial construction; either both resources are provided and a
verdict is produced, or the program does not type-check.

In the implementation, \texttt{Verdict::new} is the concrete
realisation of the $\otimes$-introduction rule
(Axiom~\ref{ax:intro}):

\begin{verbatim}
pub fn new(
    evidence:   EvidenceToken,
    compliance: ComplianceToken,
    decision:   Decision,
    explanation: String,
) -> Verdict
\end{verbatim}

Both \texttt{evidence} and \texttt{compliance} are taken \emph{by
value}.  Rust's ownership semantics enforce the linear discipline:
after the call, neither token is accessible to the caller; they have
been irreversibly consumed in the construction of \texttt{Verdict}.
The HMAC integrity seal computed inside \texttt{new} binds
\texttt{evidence\_id}, \texttt{decision}, and \texttt{explanation}
into a tamper-evident record, making post-hoc alteration detectable.

This atomicity property is stronger than a runtime assertion.  An
\texttt{assert(evidence != null)} can be removed by a compiler
optimisation, disabled under load, or bypassed by a code path that
was never tested.  The linear type law is enforced by the type
checker; it applies to every compilation, every code path, and every
version of the program.

% ----------------------------------------------------------------
\subsection{Encapsulation as Compile-Time Enforcement Boundary}
\label{sec:encapsulation}

The type law alone is insufficient if an adversary (or an
inattentive developer) can construct a \texttt{Verdict} by other
means---for instance, by writing a struct literal or by directly
assembling the hash value.  Three complementary protections close
the encapsulation perimeter:

\begin{table}[t]
\caption{The Three Protections of the Constitutional Enclosure}
\label{tab:protections}
\begin{tabular}{llll}
\hline
\textbf{ID} & \textbf{Attack Vector} & \textbf{Mechanism} & \textbf{Compiler Error} \\
\hline
A & Struct-literal \texttt{Verdict \{ \ldots \}} &
  All fields \texttt{private} & \texttt{E0451} \\
B & External \texttt{consume()} call &
  \texttt{pub(crate)} visibility & \texttt{E0603} \\
C & Silent drop of \texttt{EvidenceToken} &
  \texttt{\#[must\_use]} + deny & \texttt{unused\_must\_use} \\
\hline
\end{tabular}
\end{table}

Protection~A (Table~\ref{tab:protections}) prevents an external
caller from assembling a \texttt{Verdict} without going through the
constructor, regardless of whether they know the field types.
Protection~B prevents an external caller from extracting the inner
\texttt{Blake3Hash} out of an \texttt{EvidenceToken} and reusing it.
The \texttt{pub(crate)} qualifier restricts \texttt{consume()} to
the defining crate; any external call is rejected at compile time.
Protection~C prevents a developer from acknowledging an
\texttt{EvidenceToken} and then discarding it---the \texttt{\#[must\_use]}
attribute, combined with a crate-level \texttt{\#![deny(unused\_must\_use)]},
promotes the warning to a hard error.

Together, these three protections form a \emph{closed perimeter}: the
only well-typed path to a \texttt{Verdict} in any external crate is
through \texttt{Verdict::new}, which enforces the $V \multimap (E \otimes C)$
law.

\paragraph{Scope of the Theorem.}
The Constitutional Enclosure Theorem (Section~\ref{sec:theorem})
holds within the following boundaries:

\begin{itemize}
  \item \textbf{In scope:} All Safe Rust code.  The type checker
    enforces the visibility and ownership rules without exception.
  \item \textbf{Out of scope:} Code using the \texttt{unsafe}
    keyword.  \texttt{unsafe} permits raw pointer manipulation and
    can bypass visibility rules.  Deployments should enforce a
    zero-\texttt{unsafe} policy via \texttt{cargo geiger}.
  \item \textbf{Out of scope:} Semantic truthfulness of the evidence.
    The theorem guarantees that \emph{some} context was hashed, not
    that the context is the correct or complete input.  Policy-level
    correctness is a separate concern addressed at the application layer.
  \item \textbf{Out of scope:} Policy accuracy.  Whether the
    \texttt{ComplianceToken} accurately reflects the applicable law
    is a legal question, not a type question.
\end{itemize}

The boundary conditions above are not weaknesses of the approach;
they reflect a deliberate separation of concerns.  The type system
enforces \emph{structural} accountability---that a decision cannot
exist without evidence.  \emph{Substantive} accountability---that
the evidence is accurate and the policy is correct---is the province
of audit, certification, and human review.
